\begin{frame}{Trabalho Relacionado}
    \begin{columns}[T]
        % Coluna para Contexto
        \column{0.48\textwidth}
        \begin{block}{Contexto}
            \begin{itemize}
                \item \cite{brasnam2024F1} estudam a \textit{influência} de pilotos de Fórmula 1.
                \item Foco em \textit{pódios} como métrica principal.
                \item Dados: 855 pilotos, 1,079 corridas.
                \item Métodos: análise de redes com centralidade e \textit{PageRank}.
            \end{itemize}
        \end{block}
        
        % Coluna para Resultados
        \column{0.48\textwidth}
        \begin{block}{Resultados}
            \begin{itemize}
                \item \textit{PageRank} destaca pilotos influentes por pódios.
                \item Supera métricas tradicionais (ex.: vitórias).
                \item Propõe um novo \textit{padrão de excelência}.
                \item Aplicável à Fórmula 1 e outros esportes.
            \end{itemize}
        \end{block}
    \end{columns}
    
    % Referência
    \vspace{0.3cm}
    \small{Faria, G. et al. "Podium and Influence: A Network Analysis of the Most Important Formula One Drivers." \textit{Proceedings of the Brazilian Symposium on Multimedia and the Web}, 2021, pp. 217--224, ACM.}
\end{frame}

\begin{frame}{Trabalho Relacionado}
    \begin{columns}[T]
        % Coluna para Contexto
        \column{0.48\textwidth}
        \begin{block}{Contexto}
            \begin{itemize}
                \item \cite{buldu2019complex} analisam o \textit{desempenho} de nadadores em Jogos Olímpicos.
 Foco em \textit{modelos de redes complexas}.
                \item Uso de \textit{decomposição espectral} para estudar padrões.
                \item Dados: resultados de competições olímpicas.
            \end{itemize}
        \end{block}
        
        % Coluna para Resultados
        \column{0.48\textwidth}
        \begin{block}{Resultados}
            \begin{itemize}
                \item Revela \textit{padrões de desempenho} entre nadadores.
                \item Identifica conexões em redes de performance.
                \item Aplicável a outros esportes competitivos.
            \end{itemize}
        \end{block}
    \end{columns}
    
    % Referência
    \vspace{0.3cm}
    \small{Pereira, F. et al. "Complex Networks Models and Spectral Decomposition in the Analysis of Swimming Athletes’ Performance at Olympic Games." \textit{Frontiers in Physiology}, vol. 10, 2019, p. 1134, Frontiers, doi:10.3389/fphys.2019.01134.}
\end{frame}